\subsection{Problem Setup}
We consider post-generation answer selection under a fixed per-method budget
parameter (budget $=6$) on Cohere command-r-plus-08-2024 runs over
\texttt{openai/gsm8k} examples. GSM8K, a grade-school math benchmark, was
introduced by \citet{cobbe2021training} and provides 1,319 test problems with
numerical final answers that are amenable to exact-match evaluation. For each
example, candidate-producing methods produce final-answer candidates under the
same provider, model, and budget cap, together with method-specific runtime
metadata. A downstream selector then chooses a final answer from those
already-generated candidates. The objective is to maximize exact-match accuracy
without using gold labels at runtime.

This framing separates \emph{candidate generation} from \emph{answer selection}.
Candidate producers are the methods that spend the matched budget to generate
answers. The selector is a post-generation policy that observes candidate answers
and runtime-visible metadata and may either keep the default frontier answer or
defer to another answer source. No new reasoning branches are sampled during this
selection step.

The frontier method is a direct-reserve semantic frontier generator. At a high
level, it allocates the fixed budget across a small set of direct-reserve attempts
and frontier exploration branches, canonicalizes their final answers into answer
groups, and selects a default frontier answer from the resulting candidate pool.
During this execution it records runtime trace metadata: answer-group counts,
support for the selected frontier group, whether the direct-reserve answer agrees
with the frontier answer, and an \texttt{override\_reason} field summarizing why
the frontier selector committed to its final answer. The selector proposed in this
paper treats those metadata as diagnostics, not as correctness labels.

Two \texttt{override\_reason} values are central. The value
\texttt{single\_weak\_frontier\_branch} is emitted when the selected frontier
answer is supported by a narrow or weak frontier trace, which suggests limited
exploration of alternative answer groups. The value
\texttt{direct\_frontier\_agree} is emitted when the direct-reserve and frontier
answers agree internally. Both fields are produced before any comparison to the
gold answer; they are therefore runtime-visible and gold-free.

The selective-deferral view of the problem is therefore: given a candidate pool
and associated runtime-observable features, identify a subset of examples on which
retaining the frontier answer is less reliable than deferring to another answer
source. In this manuscript, the alternative sources are the three external
baselines L1, S1, and TALE, and the final selector is deterministic.

Evaluation is reported on:
\begin{itemize}
    \item final300 (seed 71),
    \item aggregate720 disjoint primary set (seeds 41, 61, 71),
    \item optional sensitivity set (+100; labeled sensitivity-only).
\end{itemize}

\subsection{Terminology}
\label{sec:terminology}

The following terms are used consistently throughout the paper.
Where helpful for readability, we use short manuscript names and keep the corresponding implementation identifiers in Table~\ref{tab:name_mapping} for reproducibility.

\begin{description}[leftmargin=2.2cm, labelwidth=2cm, labelsep=0.2cm, style=nextline]
  \item[Candidate pool] The set of final answers produced by the available methods (frontier and external baselines) for a given example. Candidate-pool construction is separate from the downstream selection policy.
  \item[Frontier method] The primary, computationally intensive answer selected by the direct-reserve semantic frontier policy (\texttt{direct\_reserve\_semantic\_frontier\_v2}), which serves as the default output before any policy gate fires.
  \item[External baseline] Any of the three independently-sourced single-model methods used as reference signals: L1, S1, or TALE (defined below). External baselines provide gold-free consensus signals without access to gold labels.
  \item[L1] An external method that selects the answer with the highest support count across sampled model outputs.
  \item[S1] An external method that applies sequential budget-forcing to elicit a final answer within the allocated budget.
  \item[TALE] An external method that uses prompt-level budget cues to guide answer selection.
  \item[Logical API call] A single provider inference request counted against a candidate-producing method's budget.
  \item[Budget] The integer upper bound $B = 6$ on logical API calls per example for each candidate-producing method, applied identically to all compared methods (matched-budget condition).
  \item[Matched budget] All compared candidate-producing methods operate under the same per-method budget cap $B$, ensuring that accuracy differences are not confounded by unequal per-method computational limits. If a deployment generates frontier, L1, S1, and TALE outputs from scratch, its total wall-clock and API cost depends on which producers are executed.
  \item[Postrun replay] Offline re-evaluation of a previously generated run artifact using a policy rule without issuing new provider API calls. Postrun replay is used to compute policy-specific accuracy from existing logs.
  \item[Gold-free runtime feature] Any input to a policy decision rule that is derived solely from model outputs, run metadata, or structural properties of the candidate pool---never from gold labels, exact-match signals, or correctness indicators.
  \item[Held-out evaluation evidence] Evidence collected on evaluation sets that were not used during rule development, meeting the disjointness, seed-independence, and aggregate-scale criteria documented in Section~\ref{sec:experiments}.
  \item[Provider-specific evidence] Results obtained from a particular provider--model pairing (here, Cohere command-r-plus-08-2024). Claims are explicitly scoped to this setting and do not generalize beyond it without additional replication.
\end{description}
